\chapter{Conclusion and Future Work\label{chapter:conclusions}}

This thesis compared modern C++ and Rust on a UE Context Management service modelled on a real
Nokia gNB-CU-CP component, along four dimensions: tail-latency predictability, compile-time
defect prevention, safety-model overhead, and incremental-adoption cost. Built against one frozen
specification and driven through identical code paths, the two implementations isolate language
effects from design and skill.

Three conclusions follow. First, on safety the difference is one of \emph{kind}, with a clear
boundary. Safe Rust rejects all seven injected memory-safety and data-race classes at compile
time, one of them a class no C++ sanitiser reports at all, so hazards C++ can only hope to catch
under an instrumented test become unwriteable. The edge stops there. All seven boundary classes
compile in safe Rust and all seven misbehave, though not always as they do in C++, and on two of
them the C++ tooling is the one that reports. The guarantee is therefore best described precisely:
Rust eliminates data races, not race conditions, and aliasing errors, not lock-ordering errors,
leaks, wrapping arithmetic or fault containment, and it eliminates none of them in a region marked
\texttt{unsafe}. Where the boundary defects are caught, they are caught at test time, in two of
the three cases only after being rewritten into a model checker's primitives, which is a real cost
and a different one from the sanitiser's.

Second, on migration the cost is dominated by API design and the learning curve, not by any
run-time tax. An unmodified Rust core can be driven from C++ across a thin,
\texttt{unsafe}-isolated seam whose cost on a real operation, measured against the same operation
called natively, is bounded at $\FfiOpBoundPct\%$. Nor do a build engineer's first objections
change that ordering. Most of the Rust build's extra expense is compiling dependencies and not the
compiler's own work, and the binary-size gap is a linkage policy whose ordering depends on which
matched policy is chosen, Rust smaller when both standard libraries are linked statically and C++
smaller when both are resolved dynamically. The compared component resolves no third-party code at
all. Where C++ leads is the edit-compile loop, and only for edits that stay inside a translation
unit.

Third, on performance the evidence first divides by what is being measured. An end-to-end
operation here is $\BudgetCoreShareMax\%$ the work the language performs, the rest a cross-thread
handoff the operating system performs, so an end-to-end null cannot bound a language effect at
all. Measured on the operation alone, Rust is faster, by $\BudgetCoreRatio\times$, for reasons
tracing to the standard library's hash-map layout and not to anything about safety. At the
representative load the paired test resolves $\PFiftyDiffAbsMax\,\mu$s in Rust's favour under both
allocators, too small for a deployment to care about at the median, and at p99.99 the slowest Rust
variant sits $\FarTailLangSep\,\mu$s below the fastest C++ one. The Rust arms hold their offered
rate to $\TrackingLoadRust$~kops/s against $\TrackingLoadCpp$. Beyond that the C++ build sustains
$\OverloadThrRatio\times$ the Rust build's throughput, in \OverloadThrAgree{} interleaved repeats
under both allocators, while Rust spends $\OverloadCpuPerOpRatio\times$ the CPU on each operation
it completes, so what it loses is not idleness but effort that finishes nothing. The isolated
primitives do not reproduce that ordering at significance (Section~\ref{sec:contention}). On
memory Rust carries a few MiB more fixed footprint under both allocators, in the same direction
and order of magnitude on every machine measured, while the cost of the allocator choice moved by
more than tenfold between them.

Taken together these support a narrow claim. For stateful, concurrent, partitionable
control-plane services of this kind, at the loads where the service still meets its offered rate,
the choice is which classes of defect an organisation wants to make unwriteable, and whether it
will pay the engineering and learning cost of getting them. Past that point it is a performance
decision as well, for reasons lying in the concurrency plumbing and not in either language's
guarantees. What this thesis adds is an instrument and not a verdict: the frozen specification,
the parity-tested pair, a measurement framework reaching past the saturation knee, a decomposition
that says how much of what it measures the language is answerable for, and a defect catalogue that
maps where the guarantee stops.

Several directions extend the work. The immediate priority needs no new machine: vary the
transparent-hugepage policy on one host with everything else held fixed. Three hosts agree that
the Rust build carries the larger footprint and disagree about how that compares against the
allocator cost (Section~\ref{sec:hostcompare}), the hugepage policy differed between them, and
only the system-allocator arm inherits it. One machine and two settings settle whether a footprint
claim of this kind belongs to an allocator or to a kernel setting. Beyond that, the framework can
take captured anonymised traces in place of the synthetic workload, further control-plane
subsystems or the data-heavy L2/PHY path, and a whole-program model checker covering foreign
dependencies~\cite{pearce2026rustmc}. It can also compare Rust async runtimes against
thread-per-shard, separating the tail divergence Arvidsson and Kaller attribute to
\texttt{async}/\texttt{await}~\cite{arvidsson2025async}, measure in-process isolation of the two
heaps at the seam~\cite{rivera2021galeed,kirth2022pkrusafe}, and track defect density and review
effort as Rust is adopted in a production RAN codebase.

% Declaration of generative-AI use, required by the University when a language model has
% been used. Shared with Thesis_v2.tex; edit ai_statement.tex, not a copy of it.
%% Declaration of generative-AI use. The University of Helsinki requires a student who
%% has used a language model in preparing a thesis to state in writing which model was
%% used and in what way; the model may not be named as an author or cited as a source.
%%
%% Shared verbatim by Thesis_POC.tex and Thesis_v1.tex so the two cannot drift apart.
%% It uses no generated macros and no chapter numbers, because the two documents number
%% their chapters differently.
%%
%% The order of the paragraphs follows the University's basic reporting instruction:
%% what the use covered, what it was for, which tool, how it performed and how its
%% output was checked, and what the author took from it.

\cleardoublepage
\phantomsection
\addcontentsline{toc}{chapter}{Statement on the Use of AI Tools}
\chapter*{Statement on the Use of AI Tools}
\markboth{Statement on the Use of AI Tools}{Statement on the Use of AI Tools}

\noindent\textbf{Tool.} Anthropic's Claude Opus~5, used through the Cursor development
environment, is the only generative AI system used in the preparation of this thesis. It was
used throughout the work, from June to August~2026. Following University of Helsinki policy the
model is neither named as an author nor cited as a source: it cannot take responsibility for
what it produces, and it cannot reliably vouch for the literature it paraphrases.

\noindent\textbf{What the assistance covered.} It was extensive, and is declared as such. The
model drafted and revised the running prose of this thesis and proposed the order in which its
arguments are made. It wrote and refactored the two implementations of the UE Context
Management service, the load generator, the fault-injection catalogue, and the interoperability
prototype. It wrote the analysis code behind the bootstrap confidence intervals, the exact
paired permutation tests, and the Holm correction, together with the figure-generating and
build scripts and the \LaTeX{} sources. It also suggested candidate literature.
% ==================== ACTION REQUIRED BEFORE SUBMISSION ==========================
% The abstract of Thesis_POC.tex was drafted with model assistance. The University's
% thesis guide states that a language model may not be used when preparing the thesis
% summary, and in this programme the summary also serves as the maturity test, where
% AI use is never permitted. Rewrite the abstract unaided before submitting; until
% that is done the closing sentence of this paragraph is not yet true of Thesis_POC.
% =================================================================================
The abstract was written by the author unaided, as the University's thesis guide requires of
the thesis summary.

\noindent\textbf{What the model did not settle.} The research problem, the four research
questions, and the three hypotheses are the author's, as is the control on which the study's
value rests: freezing one language-neutral specification and one concurrency model, so that the
language remains the only systematic difference between the two implementations. So is the
decision that keeps RQ2 falsifiable, namely to carry two defect classes -- a lock-order deadlock
and a lost update -- that lie outside Rust's compile-time guarantee and can therefore embarrass
the premise rather than confirm it. Every measurement reported here was run by the author on the
hosts documented in this thesis; no measured value came from a language model. Each measured
number reaches the text through a macro regenerated from the raw logs, so a figure cannot be
transcribed by hand, by the author or by the model. The readings of those numbers are the
author's as well, including the two restraints the evidence imposes: that a far tail a
virtualised host cannot resolve must not be reported as equivalence, and that a sweep dominated
by cross-thread handoff cannot bound an effect of the language.

\noindent\textbf{How the output was checked.} Model output was treated as a draft to be
falsified rather than as a result. Claims were checked against the sources cited for them: every
bibliography entry was resolved against Crossref, DataCite, or the issuing body, and each was
traced to the page or clause carrying the claim made of it. That audit found four passages where
the prose overstated or misplaced its source, one of them a sentence inside quotation marks that
does not appear in the work it was attributed to; all four were corrected, and the audit is kept
with the artefacts so that the correction remains reviewable. Code was checked by running it:
the two implementations are held to one conformance suite whose transcripts are diffed line by
line, and the fault catalogue fixes the expected outcome of each defect before it is injected,
so that a surprise appears as a discrepancy instead of being explained away afterwards. Nothing
is asserted here on the model's authority.

\noindent\textbf{What the author took from it.} The model is dependable at mechanical work --
transposing an idiom between two languages, restructuring a table, keeping a build honest -- and
undependable in a specific and predictable way: it will supply a plausible citation for a claim,
a plausible cause for a measurement, and a confident sentence where the data supports a hedge.
Its fluency is not evidence. The lasting lesson of the work was accordingly less about either
language than about verification. The effort saved in drafting is spent again, and rightly, in
checking, and the parts of this thesis that are defensible are the parts that were checked
against a source, a specification, or a measurement.

\noindent\textbf{Responsibility.} The author is responsible for the whole of this thesis: its
language, its facts, its code, its measurements, and its conclusions. Where it is wrong, the
error is the author's.


\cleardoublepage %fixes the position of bibliography in bookmarks
\phantomsection

\addcontentsline{toc}{chapter}{\bibname}
\printbibliography
